\chapter{SDV Integration with Eclipse Leda}\label{ch:sdv}
\thispagestyle{pagestyle}

\section{Eclipse Leda Overview}
\label{sec:leda-overview}

Eclipse Leda \cite{EclipseLeda2023} is an SDV distribution built on the
Yocto Project. It packages a curated set of Eclipse SDV middleware
components --- Mosquitto MQTT broker, Kuksa databroker, Kanto container
manager, and OpenTelemetry collector --- into a bootable image that runs
on QEMU (x86-64) or on supported ARM hardware such as the Raspberry Pi~4.
Its primary purpose is to serve as a reference platform for developing
and validating SDV software stacks before deploying to a real vehicle
ECU.

SHILATE targets the QEMU image
(\texttt{sdv-image-all-qemux86-64.wic.qcow2}) included in the
repository. The \texttt{leda/run-leda.sh} script launches the image with
a bridged network interface so that the host machine and the QEMU guest
share a Mosquitto broker at \texttt{127.0.0.1:1883}.

\section{MQTT-to-Kuksa Bridge}
\label{sec:feeder}

The \texttt{mqtt-kuksa-feeder} container subscribes to the
\texttt{vehicle/\#} wildcard on the Mosquitto broker and forwards each
received message to the Kuksa databroker \cite{EclipseKuksa2024} as a
VSS path write via gRPC. The topic-to-VSS mapping is shown in
\cref{table:vss-map}.

\begin{table}[H]
\caption{MQTT topic to VSS path mapping.}
\label{table:vss-map}
\begin{tabular}{|p{5cm}|p{7cm}|}
  \hline
  \textbf{MQTT topic} & \textbf{VSS path} \\
  \hline
  \texttt{vehicle/speed}    & \texttt{Vehicle.Speed} \\
  \hline
  \texttt{vehicle/rpm}      & \texttt{Vehicle.Powertrain.CombustionEngine.Speed} \\
  \hline
  \texttt{vehicle/steering} & \texttt{Vehicle.Chassis.SteeringWheel.Angle} \\
  \hline
  \texttt{vehicle/brake}    & \texttt{Vehicle.Chassis.Brake.PedalPosition} \\
  \hline
  \texttt{vehicle/throttle} & \texttt{Vehicle.OBD.ThrottlePosition} \\
  \hline
\end{tabular}
\end{table}

After the feeder writes a signal, any subscriber --- including the
Velocitas vehicle application --- is notified with the new value within
the same event loop tick.

\section{Velocitas Vehicle Application}
\label{sec:velocitas}

The Velocitas vehicle application \cite{EclipseVelocitas2024}
implements the operational safety monitor for the RL agent. It subscribes
to five VSS signals via the Kuksa databroker, logs a telemetry summary
every 5 seconds, and enforces two hard limits:
\begin{itemize}[topsep=1.15pt,itemsep=1.15pt]
  \item \textbf{Overspeed}: if \texttt{Vehicle.Speed} exceeds 120~km/h,
    an alert is published to \texttt{leda/command/alert} with payload
    \texttt{\{"type":"overspeed","value":<speed>\}}.
  \item \textbf{RPM redline}: if engine RPM exceeds 6500, an analogous
    alert is published.
\end{itemize}
Alerts are currently informational; the architecture is designed so that
the safety monitor can be upgraded to a command-injection path without
modifying any other component.

\section{Container Deployment}
\label{sec:containers}

All three service containers (\texttt{leda-controller},
\texttt{mqtt-kuksa-feeder}, \texttt{velocitas-app}) ship with Dockerfiles
producing minimal Python~3.11 slim images. Each is described by a Kanto
container manifest (\texttt{kanto-manifest.json}) specifying the OCI
image reference, environment variables
(\texttt{MQTT\_HOST}, \texttt{MQTT\_PORT}, \texttt{MODE},
\texttt{MODEL\_PATH}), and restart policy. Provisioning is automated by
\texttt{leda/provision-device.sh}, which uses SSH to copy manifests to
the Leda device and invoke the Kanto CLI to pull and start the
containers.

\section{Leda Telemetry Broker}
\label{sec:telemetry-broker}

In addition to MQTT, the Leda image runs a lightweight telemetry broker
implemented in Rust. The broker listens on TCP port~7878 and acts as a
multiplexing relay: any client that connects and sends telemetry data has
that data forwarded to all other connected clients. This enables tools
such as \texttt{netcat} or custom dashboards to observe the live
telemetry stream without subscribing to individual MQTT topics.

\section{Raspberry Pi 4 Hardware Actuator}
\label{sec:rpi-actuator}

Closing the loop between the SDV software layer and a real physical
device requires hardware that can translate a VSS signal into a
mechanical action. The design described in this section is theoretical
--- the physical build has not been assembled --- but it is grounded in
components that are readily available and well understood.

The target application is a 12~V DC water pump driven by a
pulse-width-modulated (PWM) signal, where the duty cycle is set by the
RL agent's throttle output via the VSS path
\texttt{Vehicle.OBD.ThrottlePosition}.

\subsection{Circuit Design}
\label{subsec:circuit}

A Raspberry Pi 4B \cite{RaspberryPi4} runs Eclipse Leda on its ARM
processor and therefore participates in the same SDV software stack as
the QEMU simulation. Its 3.3~V GPIO pins cannot drive a 12~V motor load
directly, so a MOSFET module is interposed. A suitable choice is the
IRF520N N-channel MOSFET: its gate threshold voltage is well within the
3.3~V logic range, its $R_{\mathrm{DS(on)}}$ is low enough
(approx.\ 0.27~$\Omega$) to avoid significant dissipation at pump
currents, and its $V_{\mathrm{DS}}$ rating of 100~V comfortably exceeds
the 12~V supply.

The gate is driven through a 1~k$\Omega$ series resistor to limit
inrush current and damp potential ringing on the gate trace. A small
pull-down resistor (10~k$\Omega$) from gate to ground ensures the
MOSFET is fully off when the GPIO pin is tri-stated at boot, preventing
unintended motor activation during start-up.

The pump motor is an inductive load. When the MOSFET switches off
abruptly, the energy stored in the motor's inductance $L$ has nowhere
to go and appears as a voltage spike across the switching node:
\begin{equation}
  V_{\mathrm{spike}} = -L \frac{dI}{dt}.
  \label{eq:inductive-kick}
\end{equation}
Without protection, this spike can exceed the MOSFET's $V_{\mathrm{DS}}$
rating and destroy the transistor. A \emph{flyback suppressor} --- also
called a freewheeling or snubber diode --- provides a low-impedance
current path that clamps the spike. A 1N4007 general-purpose rectifier
connected in reverse across the pump terminals (cathode to the 12~V rail,
anode to the drain) is sufficient for modest pump currents. For faster
PWM frequencies or higher currents, a Zener diode or a dedicated TVS
(Transient Voltage Suppressor) device placed in the same position offers
a tighter voltage clamp.

The complete circuit, described schematically, is:

\begin{quote}
  \textbf{RPi GPIO23} $\xrightarrow{\text{1 k}\Omega}$ \textbf{MOSFET Gate}
  --- \textbf{MOSFET Drain} $\to$ \textbf{Pump (-)};
  \textbf{Pump (+)} $\to$ \textbf{12 V};
  \textbf{1N4007/TVS} from Drain to 12 V (reverse-biased);
  \textbf{MOSFET Source} $\to$ \textbf{GND};
  \textbf{10 k}\boldmath$\Omega$ \textbf{pull-down} from Gate to GND.
\end{quote}

\subsection{Software: VSS Subscription to PWM Output}
\label{subsec:rpi-software}

The Raspberry Pi runs a small Python daemon alongside the Leda
containers. It subscribes to the Kuksa databroker using the
\texttt{kuksa-client} library and, on each value change,
converts the normalised throttle reading to a PWM duty cycle for
GPIO~23:

\begin{lstlisting}[language=Python,
  caption={Theoretical Raspberry Pi 4 pump controller subscribed to VSS via Kuksa.},
  label={code:rpi}]
import asyncio
import RPi.GPIO as GPIO
from kuksa_client.grpc.aio import VSSClient

PUMP_GPIO  = 23
PWM_FREQ   = 1000   # Hz -- 1 kHz carrier well above audible range
VSS_PATH   = "Vehicle.OBD.ThrottlePosition"
KUKSA_ADDR = ("127.0.0.1", 55555)

GPIO.setmode(GPIO.BCM)
GPIO.setup(PUMP_GPIO, GPIO.OUT)
pwm = GPIO.PWM(PUMP_GPIO, PWM_FREQ)
pwm.start(0)

async def main() -> None:
    async with VSSClient(*KUKSA_ADDR) as client:
        async for updates in client.subscribe_current_values([VSS_PATH]):
            for path, value in updates.items():
                # ThrottlePosition is 0.0-100.0 (percent); map to duty cycle
                duty = max(0.0, min(100.0, float(value.value)))
                pwm.ChangeDutyCycle(duty)

asyncio.run(main())
\end{lstlisting}

The Kuksa gRPC subscription fires immediately on value change, so the
pump duty cycle tracks the agent's throttle output with sub-millisecond
software latency. The PWM carrier frequency of 1~kHz is above the
audible range of most motors, reducing acoustic noise without requiring
an especially low MOSFET gate capacitance.

\subsection{Connection to SHILATE}
\label{subsec:rpi-connection}

Within the SHILATE data flow, the RL agent running on the host machine
publishes a throttle action to the MQTT topic
\texttt{env0/vehicle/control/throttle}. The \texttt{mqtt-kuksa-feeder}
container forwards this to the Kuksa databroker as
\texttt{Vehicle.OBD.ThrottlePosition}. The RPi daemon described above
subscribes to that path and drives the pump accordingly.

This is precisely the Software-Defined Vehicle pattern: the physical
actuator has no knowledge of the RL algorithm, the simulation, or the
MQTT broker. It speaks only VSS signals. Replacing the pump with any
other VSS-aware actuator --- a servo, a heater, a relay --- requires no
changes to the Python training code.
